% !TEX root = ../main.tex
\chapter{绪\hspace{1em}论}[Introduction]

\section{课题背景及研究的目的和意义}[Background, objective and significance]

高保真偏微分方程求解器在反问题、最优控制与不确定性量化中仍然昂贵。数据驱动代理模型把求解过程近似为一次前向推断，使大规模参数扫描成为可能。然而，当运行工况离开训练分布时，残差不再服从训练阶段观察到的规律，名义置信区间会系统性地低估风险。

本课题的目标是：在\textbf{不重新训练}冻结代理模型的前提下，利用较小的标注校准集恢复区间覆盖率，并为明显失效的工况提供回到数值求解器的判据。

\section{国内外研究现状}[Related work]

物理信息神经网络与神经算子为科学代理提供了可微的函数类。共形推断与分位回归则从统计侧给出有限样本覆盖保证。将二者结合的工作仍然较少：多数科学机器学习论文报告点误差，而不是校准后的区间。本文把残差分数上的单调映射作为连接两者的最小接口。

\section{本文的主要研究内容}[Main research contents]

本文的工作组织如下：第2章给出残差分数、经验分布与覆盖率命题；结论章汇总创新点与后续工作。附录给出补充表格与符号说明。
